% Generic Chinese computer-science survey template.
% Replace the document class with the official target-journal template before submission.
\documentclass[UTF8,a4paper,10.5pt]{ctexart}

\usepackage[margin=2.4cm]{geometry}
\usepackage{booktabs}
\usepackage{longtable}
\usepackage{graphicx}
\usepackage{amsmath,amssymb}
\usepackage{hyperref}
\usepackage[numbers,sort&compress]{natbib}
\usepackage{enumitem}

\title{综述题目：对象、技术脉络与发展趋势}
\author{作者姓名\\单位名称}
\date{}

\begin{document}
\maketitle

\begin{abstract}
说明研究背景与重要性，明确综述范围和检索截止日期，概括本文采用的分类主线、核心比较结论、关键挑战与未来方向。摘要不应写成章节目录，也不应使用未在正文解释的缩写。
\end{abstract}

\noindent\textbf{关键词：} 关键词1；关键词2；关键词3；关键词4

\section{引言}
\subsection{研究背景与问题提出}
\subsection{已有综述及其不足}
\subsection{本文范围、方法与主要贡献}
\subsection{文章组织结构}

\section{基础概念与问题模型}
统一术语、系统模型、威胁模型、评价指标和必要的技术背景。

\section{研究分类体系}
先说明分类轴及选择理由，再逐类展开。分类应服务于比较与理解，而不是简单按论文名称分组。

\subsection{技术路线一}
\subsection{技术路线二}
\subsection{技术路线三}

\section{研究演化与关键里程碑}
从问题、方法与系统条件的变化解释阶段转折，并配合时间线或引用关系图。

\section{典型方法比较}

\begin{table}[htbp]
\centering
\caption{典型方法比较}
\begin{tabular}{lllll}
\toprule
方法 & 问题模型 & 核心机制 & 优势 & 局限 \\
\midrule
方法A &  &  &  &  \\
方法B &  &  &  &  \\
\bottomrule
\end{tabular}
\end{table}

\section{应用、数据集与评价方法}
比较应用场景、公开数据集、实验协议、指标及其可比性问题。

\section{关键挑战}
每项挑战应包含：已有证据、尚未解决的技术障碍、适用边界与评价标准。

\section{未来发展方向}
未来方向必须由前文证据推导，并说明潜在技术路径、风险以及可验证的研究问题。

\section{总结}
回答本文提出的研究问题，重申最重要的综合判断，不引入新的文献或结论。

\section*{综述局限性}
说明数据库覆盖、语言范围、时间边界、全文可得性和主观分类决策等局限。

\bibliographystyle{unsrtnat}
\bibliography{references}

\end{document}
